\def\level{BTS}  % Classe à droite
\def\course{CIEL 2}
\def\chaptername{Lois continues} % Titre au centre
\def\docpart{Loi exponentielle~--~TP3} % Partie du cours
\def\docname{C228}        % Identifiant à gauche

\pagestyle{cours_FP}
\makeheaderBTS{} % Affiche le bandeau


\begin{exercice_cours}[\quad]{}
	La durée de vie d'un robot, exprimée en années est modélisée par une variable aléatoire $X$ qui suit la loi exponentielle de paramètre $\lambda=0,2$.

	\begin{enumerate}
		\item A quel instant $t$, à un mois près, la probabilité qu'un robot tombe en panne pour la première fois est-elle de 0,5?
		\item Montrer que la probabibilité qu'un robot n'ait pas eu de panne au cours des deux premières années est $e^{-0,4}$.
		\item Calculer la probabilité qu'un robot tombe en panne au cours de la troisième année sachant qu'il n'a pas eu de panne au cours des deux premières années.
	\end{enumerate}
\end{exercice_cours}

\begin{exercice_cours}[\quad]{}

On s'intéresse à la durée de vie d'une ampoule électrique. On suppose que cette durée de vie, exprimée en heures, est modélisée par une variable aléatoire $X$ qui suit la loi exponentielle de paramètre $\lambda=0,01$.

\begin{enumerate}
	\item Calculer $P(X\leqslant 100)$, 
	\item Calculer $P(50\leqslant X\leqslant 150)$ 
	\item Calculer la probabilité que l'ampoule dure plus de 200 heures.
	\item Calculer la probabilité que l'ampoule dure exactement 100 heures.
	\item Calculer l'espérance de $X$ et donner une interprétation de ce résultat.	
	\item Calculer la probabilité que l'ampoule dure plus de 150 heures sachant qu'elle a déjà duré 100 heures.
\end{enumerate}
\end{exercice_cours}

\begin{exercice_cours}[\quad]{}

	On s'intéresse à la durée de vie d'une batterie de téléphone portable. On suppose que cette durée de vie, exprimée en heures, est modélisée par une variable aléatoire $Y$ qui suit la loi exponentielle de paramètre $\lambda$.

	\begin{enumerate}
		\item Tracer la fonction de répartition de $Y$ à l'aide de Géogébra pour différentes valeurs de $\lambda$.
		\item Sachant que $P(Y\leqslant 10)=0,6321$, déterminer de $\lambda$ grâce à Géogébra.
		\item Vérifier la valeur trouvée par le calcul.
		\item Déterminer $P(Y\leqslant 20)$ grâce à Géogébra et vérifier ce résultat par le calcul.
		\item Déterminer $P(15\leqslant Y\leqslant 25)$ grâce à Géogébra et vérifier ce résultat par le calcul.
		\item Déterminer la probabilité que la batterie dure plus de 30 heures.
		\item Déterminer la probabilité que la batterie dure exactement 20 heures.
		\item Calculer l'espérance de $Y$ et donner une interprétation de ce résultat.
		\item Calculer la probabilité que la batterie dure plus de 25 heures sachant qu'elle a déjà duré 15 heures.
		\item Calculer la probabilité que la batterie dure plus de 30 heures sachant qu'elle a déjà duré 20 heures.
		\item Calculer la probabilité que la batterie dure plus de 30 heures sachant qu'elle a déjà duré 25 heures.
	\end{enumerate}
	\end{exercice_cours}


	\begin{exercice_cours}[\quad]{}


La durée de vie, en jours, d’un composant électronique est une variable aléatoire T qui suit une loi exponentielle de paramètre $\ell$ compris entre $0$ et $0,005$

{\it On rappelle que la fonction densité de probabilité associée à une loi exponentielle est : $f(t)=\ell\e^{-\ell t}$.}\\

\begin{enumerate}
	\item À l'aide de Géogebra, déterminer le paramètre $\ell$, à $0,0001$ près, pour que P(T$\leq266$)$ \simeq 0,5$.
	
	{\it \begin{itemize}
		\item Dans le champs de saisie de Géogebra, rentrer la fonction $f$.
		\item Représenter la probabilité à l'aide de la commande \texttt{intégrale(<fonction>,<xmin>,<xmax>)}.
		\item Déplacer le curseur pour obtenir la valeur demandée (il sera nécessaire d'augmenter l'arrondi et de choisir un pas adapté.
		\end{itemize}}
	
	$\ell = $\makebox[0.1\linewidth]{\dotfill} \hfill \fbox{\bf \large Appeler le professeur}
	\item[~] Dans la suite du TP, on choisira $\ell = 0,00255$.
	
	\item Calculer la MTBF (moyenne de temps de bon fonctionnement), correspondant à la durée de vie moyenne du composant éléctronique.\\
	
	\dotfill
	
	\smallskip
	\item On s'intéresse à la probabilité que ce composant tombe en panne au cours de la première année.
	
	{\it On considère qu’une année compte 365 jours.}
	
	\begin{enumerate}
		\item Donner la valeur exacte (formule) de cette probabilité : \dotfill \\
		
		\item En donner une valeur approchée à $10^{-3}$ près : \dotfill
		
		\smallskip
	\end{enumerate}
	\item On s'intéresse à la probabilité que ce composant fonctionne encore au bout de $3$ ans.
	
	Donner la valeur exacte de cette probabilité, puis une valeur approchée à $10^{-3}$ près. \\
	
	\dotfill
	
	\smallskip
	\item À l'aide de Géogebra, déterminer le plus grand nombre de jours $n$, tel que la probabilité que ce composant fonctionne encore au bout de ces $n$ jours soit supérieure ou égale à $0,95$.\\
	
	$n =$\makebox[0.1\linewidth]{\dotfill} \hfill \fbox{\bf \large Appeler le professeur}\\
	
	Retrouver ce résultat par un calcul :\\
	
	\dotfill\\
	
	\dotfill\\
	
	\dotfill\\
	
	\dotfill\\
	
	\dotfill\\
	
	\dotfill
\end{enumerate}

\end{exercice_cours}

\begin{exercice_cours}[\quad]{}
On étudie un modèle de climatiseur d’automobile composé d’un module mécanique et d’un module
électronique.

Si un module subit une panne, il est changé.

\medskip

L'enseigne d’entretien automobile a constaté que la durée de
fonctionnement (en mois) du module électronique peut être modélisée par une variable aléatoire T qui suit une loi exponentielle de paramètre $\lambda$.

\begin{enumerate}
	\item Le service statistique indique que P(T$\leq 24$) $= 0,03$.
	
	\begin{enumerate}
		\item À l'aide de Géogebra, déterminer une valeur approchée de $\lambda$ à $10^{-4}$ près.
		
		{\it Reprendre la méthode du précédent TP. On utilisera ici un curseur compris entre $0$ et $0,002$}\\
		
		$\lambda =$ \makebox[0.2\linewidth]{\dotfill} \hfill \fbox{\large \bf Appeler le professeur}
		
		\smallskip
		\item À l'aide d'un calcul, déterminer la valeur exacte de $\lambda$.\\
		
		\dotfill \\
		
		\dotfill \\
		
		\dotfill \\
		
		\dotfill \\
		
		\dotfill
		
		\smallskip
	\end{enumerate}

	
	\item[~] Dans la suite du TP, on prendra $\lambda = 0,00127$.
	
	Pour tous les résultats, on attend la valeur exacte, puis une valeur approchée à $0,01$ près.
	
	\smallskip
	
	\item Déterminer la probabilité que la durée de fonctionnement du module électronique soit comprise entre 24 et 48 mois.\\
	
	\dotfill
	
	\smallskip
	\item Déterminer la probabilité qu'un module électronique fonctionne encore au bout de 3 ans.\\
	
	\dotfill \\
	
	\dotfill
	
	\smallskip
	
	\item Le module électronique du climatiseur fonctionne depuis 36 mois.
	
	Déterminer la probabilité qu’il fonctionne encore les 12 mois suivants.\\
	
	\dotfill \\
	
	\dotfill
\end{enumerate}

\newpage

\subsection*{Exercice 2 : \normalsize  Loi exponentielle -- \small Donner les valeurs exactes puis leur valeur approchée au centième.}

La durée de vie, en années, d’un composant électronique fabriqué dans une usine est une variable aléatoire T qui suit la loi exponentielle de paramètre $\lambda = 0,099$.

\begin{enumerate}
	\item On choisit au hasard un composant fabriqué dans cette usine.
	
	Déterminer la probabilité que ce composant tombe en panne au cours des 5 premières années.\\
	
	\dotfill
	
	\smallskip
	\item On choisit au hasard un composant parmi ceux qui fonctionnent encore au bout de 2 ans.
	
	Déterminer la probabilité que ce composant ait une durée de vie supérieure à 7 ans.\\
	
	\dotfill
	
	\smallskip
	\item Donner l’espérance mathématique E(T) de la variable aléatoire T à l’unité près.
	
	Interpréter ce résultat dans le contexte de l'énoncé.\\
	
	\dotfill \\
	
	\dotfill \\
	
	\dotfill

\end{exercice_cours}